\raggedbottom

\section{Ejercicio 1: Interrupciones con botones de baja y alta prioridad
}
Para este programa se usará el código del TP1 como base, el cual se modificará para hacer interrupciones al led del TP1 con botones de baja y alta prioridad.
Para conseguir esto se agregaron las siguientes definiciones en el archivo main.c:

\begin{itemize}
    \item Se definió la dirección base \textbf{AFIO\_BASE} 
    \item Se definió la dirección base \textbf{EXTI\_BASE}
    \item Se definió la dirección base \textbf{NVIC\_BASE}
    \item Se definió el registro \textbf{AFIO\_EXTICR1}
    \item Se definió el registro \textbf{EXTI\_IMR}
    \item Se definió el registro \textbf{EXTI\_FTSR}
    \item Se definió el registro \textbf{EXTI\_PR}
    \item Se definió el registro \textbf{NVIC\_ISER0}
    \item Se definió el registro \textbf{NVIC\_IPR}
    \item Se definieron las mascaras \textbf{GPIOA6} y \textbf{GPIOA7}
\end{itemize}


\begin{figure}[H]
    \centering
    \includegraphics[width=0.49\textwidth]{images/definiciones.png}
    \caption*{Definiciones de los registros en main.c}
    \label{fig:led}
\end{figure}

Aquí se configuran los registros necesarios para poder realizar las interrupciones con los botones.
El registro \textbf{AFIO\_EXTICR1} se utiliza para configurar el pin del botón como fuente de interrupción externa; en este caso usaremos dos, el A1 y el A2.
El registro \textbf{EXTI\_IMR} se utiliza para habilitar las interrupciones de los pines configurados.
El registro \textbf{EXTI\_FTSR} se utiliza para configurar la interrupción por flanco de bajada.
El registro \textbf{NVIC\_ISER0} se utiliza para habilitar las interrupciones en el NVIC, y el registro \textbf{NVIC\_IPR} se utiliza para configurar la prioridad de las interrupciones.\\

A continuación agregamos las funciones de interrupción para los botones, las cuales se encargan de encender y apagar el led correspondiente al botón presionado. Para el botón de baja prioridad (A1) se encenderá el led del pin A6, y para el botón de alta prioridad (A2) se encenderá el led del pin A7.

\begin{figure}[H]
    \centering
    \includegraphics[width=0.4\textwidth]{images/funciones.png}
    \caption*{Agregado de funciones para interrupciones}
    \label{fig:main1}
\end{figure}

\begin{figure}[H]
    \centering
    \includegraphics[width=0.4\textwidth]{images/funciones2.png}
    \caption*{Agregado de funciones para interrupciones}
    \label{fig:main1}
\end{figure}

Ahora pasamos a modificar la función \textbf{main}

\begin{figure}[H]
    \centering
    \includegraphics[width=0.45\textwidth]{images/main.png}
    \caption*{Configuración de los pines y las interrupciones en main}
    \label{fig:main1}
\end{figure}

Primero habilitamos los clocks del GPIO y AFIO, luego configuramos los pines A5, A6 y A7 como salida, y los A1 y A2 como entrada con pull-up. Luego asociamos el pin GPIO a la línea EXTI y pasamos a configurar las interrupciones como flancos descendentes, por ultimo configuramos el NVIC con las prioridades correspondientes. 

\subsection{Imagenes del Funcionamiento}
